\section{Prospective protocols for the five discriminating tests}\label{sec:minds-protocols}

The first test should establish information gain. Recruit targets and observers, measure baseline opportunities and abilities, and randomly assign matched observation time to repeated presentation in one domain, diverse contexts, independent observers, or favorable endorsement. Define later criteria before collecting the evaluations: held-out task performance, verified follow-through, recipient benefit, or a preregistered pair outcome. Compare prediction beyond baseline information. If diverse observation raises confidence and liking without improving criterion prediction, the proposed reveal mechanism is weakened even if social preference changes.

The second test separates contribution from reputation. Randomize opportunities to perform useful work where possible, measure what recipients actually receive, and distinguish that benefit from observers' judgments. An intervention might increase support because it produces an immediate benefit, because it reveals a stable quality, or because an institution awards access. Following these routes separately prevents a favorable social outcome from being counted as evidence for all three. Measures of resources, training, health, prior ties, and opportunity should accompany the focal performance.

The third test examines directed investment and compatibility. Observe whether capability predicts contribution to this partner, rather than to an unspecified audience. Repeated dyadic interactions can test whether reliable investment changes the value placed on capability, as the main paper's directed-investment model suggests. A competing explanation is that a general positive impression drives both judgments. Independent records of delivery, multiple informants, and outcomes obtained after the evaluation help distinguish that explanation. Pair-specific forecasts should be compared against predictions based only on each person's average desirability.

The fourth test addresses authority. Record whose assessment is used, who controls the opportunity, who can refuse, and what it costs to leave. A study can compare information changes with changes in access rules, rather than assuming that better assessment necessarily produces freer choice. Family approval, peer endorsement, institutional requirements, and coercion should be measured separately. A positive association between community knowledge and a marriage outcome could arise from practical support or constrained options as well as more accurate forecasts.

The fifth test follows the outcome chain. Later contact, reciprocal partnership, relationship quality, parenthood, births, and child welfare should be recorded as distinct outcomes with appropriate time horizons. Attrition and nonresponse need explicit accounting because relationship dissolution and leaving a platform are themselves related to the processes under study. Analyses should retain people who decline opportunities or never form a partnership. Restricting follow-up to successful couples would condition on the very transition the mechanism seeks to explain.

